\documentclass{article}
\usepackage{amsmath,amssymb,amsthm}
\usepackage{algorithm}
\usepackage[noend]{algpseudocode}
\usepackage{hyperref}
\usepackage{enumitem}
\newtheorem{theorem}{Theorem}
\newtheorem{lemma}{Lemma}
\newtheorem{assumption}{Assumption}
\newcommand{\E}{\mathbb{E}}
\newcommand{\R}{\mathbb{R}}

\begin{document}

\section*{RMSProp with Running Average of Squared Gradients}

\section*{Mathematical Formulation}
Let $f:\R^d\rightarrow\R$ be a (possibly non‑convex) differentiable loss.  
At iteration $t\ge 1$ we have a stochastic gradient $g_t:=\nabla f(w_t;\xi_t)$,
where $\xi_t$ denotes the random data sample (or minibatch) drawn at step $t$.
The algorithm maintains an exponential moving average of the element‑wise
squared gradients:
\[
v_t \;=\; \beta\, v_{t-1} \;+\; (1-\beta)\, g_t \odot g_t,
\qquad v_0 = \mathbf{0},
\tag{1}
\]
where $\beta\in[0,1)$ and $\odot$ denotes element‑wise product.
The learning rate for each coordinate is adapted by the inverse square root of
the accumulated second moment:
\[
w_{t+1}\;=\; w_t \;-\; \alpha \, \frac{g_t}{\sqrt{v_t+\varepsilon}},
\tag{2}
\]
with a global stepsize $\alpha>0$ and a small constant $\varepsilon>0$ for
numerical stability.  The division and square root are taken element‑wise.

\section*{Pseudocode}
\begin{algorithm}[H]
\caption{RMSProp (scalar stepsize $\alpha$, decay $\beta$, $\varepsilon$)}
\begin{algorithmic}[1]
\State \textbf{Input:} initial point $w_0\in\R^d$, stepsize $\alpha>0$, decay $\beta\in[0,1)$, $\varepsilon>0$, max iter $T$.
\State $v_0 \gets \mathbf{0}\in\R^d$
\For{$t=0,\dots,T-1$}
    \State Sample a mini‑batch $\xi_t$ and compute stochastic gradient $g_t = \nabla f(w_t;\xi_t)$
    \State $v_{t+1} \gets \beta\, v_{t} + (1-\beta)\, g_t\odot g_t$
    \State $w_{t+1} \gets w_t - \alpha\,\dfrac{g_t}{\sqrt{v_{t+1}+\varepsilon}}$
\EndFor
\State \textbf{Output:} $w_T$ (or optionally the average $\bar w_T=\frac1T\sum_{t=1}^T w_t$)
\end{algorithmic}
\end{algorithm}

\section*{Dry Run Example}
Consider the one‑dimensional quadratic $f(w)=(w-3)^2$, whose gradient is
$\nabla f(w)=2(w-3)$.  Set $w_0=0$, $\alpha=0.1$, $\beta=0.9$, $\varepsilon=10^{-8}$.
We perform three iterations (deterministic gradients, i.e. $g_t=\nabla f(w_t)$).

\begin{enumerate}[label=Iteration \arabic*:, leftmargin=*, nosep]
\item $t=0$:
\[
g_0 = 2(0-3) = -6,\qquad
v_1 = 0.9\cdot 0 + 0.1\cdot (-6)^2 = 0.1\cdot 36 = 3.6.
\]
\[
w_1 = 0 - 0.1\frac{-6}{\sqrt{3.6+\varepsilon}}
      = 0 + 0.1\frac{6}{\sqrt{3.6}}
      \approx 0 + 0.1\frac{6}{1.897}=0.316.
\]
Objective: $f(w_1)=(0.316-3)^2\approx7.22$.

\item $t=1$:
\[
g_1 = 2(0.316-3) = -5.368,\qquad
v_2 = 0.9\cdot 3.6 + 0.1\cdot (-5.368)^2
     = 3.24 + 0.1\cdot 28.82 \approx 6.122.
\]
\[
w_2 = 0.316 - 0.1\frac{-5.368}{\sqrt{6.122}}
     = 0.316 + 0.1\frac{5.368}{2.474}
     \approx 0.316 + 0.217 = 0.533.
\]
Objective: $f(w_2)=(0.533-3)^2\approx6.09$.

\item $t=2$:
\[
g_2 = 2(0.533-3) = -4.934,\qquad
v_3 = 0.9\cdot 6.122 + 0.1\cdot (-4.934)^2
     = 5.510 + 0.1\cdot 24.34 \approx 7.944.
\]
\[
w_3 = 0.533 - 0.1\frac{-4.934}{\sqrt{7.944}}
     = 0.533 + 0.1\frac{4.934}{2.820}
     \approx 0.533 + 0.175 = 0.708.
\]
Objective: $f(w_3)=(0.708-3)^2\approx5.24$.
\end{enumerate}
The iterates move toward the minimiser $w^\star=3$, while the denominator
$\sqrt{v_t}$ grows, automatically decreasing the effective stepsize.

\newpage
\section*{Assumptions}
\begin{assumption}[Smoothness]\label{ass:smooth}
$f:\R^d\!\to\!\R$ is $L$‑smooth, i.e.
\[
\|\nabla f(x)-\nabla f(y)\|\le L\|x-y\|,\qquad\forall x,y\in\R^d.
\tag{A1}
\]
\end{assumption}
\begin{assumption}[Unbiased Stochastic Gradient]\label{ass:unbiased}
For every $t$, 
\[
\E_{\xi_t}[g_t\mid w_t]=\nabla f(w_t),\qquad 
\E_{\xi_t}[\|g_t\|^2\mid w_t]\le G^2,
\tag{A2}
\]
for some $G>0$.
\end{assumption}
\begin{assumption}[Bounded Second Moment]\label{ass:bound-second}
The variance of each coordinate is uniformly bounded:
\[
\E_{\xi_t}[g_{t,i}^2\mid w_t]\le \sigma_i^2\le \sigma_{\max}^2,\qquad
i=1,\dots,d,
\tag{A3}
\]
with $\sigma_{\max}<\infty$.
\end{assumption}
\begin{assumption}[Hyper‑parameter Range]\label{ass:hyper}
The decay $\beta\in[0,1)$ and stepsize $\alpha>0$ satisfy
\[
\alpha\le \frac{1}{L\,(1-\beta)}.
\tag{A4}
\]
\end{assumption}

\section*{Main Convergence Theorem}
\begin{theorem}[Non‑convex RMSProp convergence]\label{thm:main}
Under Assumptions \ref{ass:smooth}–\ref{ass:hyper} and for any integer $T\ge1$,
the iterates generated by \eqref{1}–\eqref{2} satisfy
\[
\frac{1}{T}\sum_{t=0}^{T-1}\E\!\left[\big\|\nabla f(w_t)\big\|^2\right]
\;\le\;
\frac{2\big(f(w_0)-f^\star\big)}{\alpha\sqrt{\varepsilon}}
\;+\;
\frac{2L\alpha G^2}{(1-\beta)\sqrt{\varepsilon}}\,\frac{\log\!\big(1+T(1-\beta)G^2/\varepsilon\big)}{T},
\tag{3}
\]
where $f^\star=\inf_{w}f(w)$. Consequently,
\[
\min_{0\le t<T}\E\!\big[\|\nabla f(w_t)\|^2\big]=O\!\left(\frac{\log T}{\sqrt{T}}\right).
\]
\end{theorem}

\section*{Theoretical Properties}
\begin{enumerate}[label=\arabic*.]
\item \textbf{Stability.} The denominator $\sqrt{v_t+\varepsilon}$ is lower bounded by
$\sqrt{\varepsilon}$, guaranteeing that the effective stepsize never exceeds
$\alpha/\sqrt{\varepsilon}$. The exponential averaging prevents $v_t$ from
exploding, which yields a bounded update magnitude.

\item \textbf{Hyper‑parameter Sensitivity.}
\begin{itemize}
\item Larger $\beta$ yields slower adaptation (more smoothing) and a larger
constant factor $(1-\beta)^{-1}$ in \eqref{3}.
\item The stepsize $\alpha$ appears linearly in the bound; choosing $\alpha$
according to \eqref{A4} ensures the descent term dominates the curvature term.
\item $\varepsilon$ trades off bias (large $\varepsilon$ reduces adaptivity) and
variance (small $\varepsilon$ may cause large stepsizes early on). The bound
depends on $1/\sqrt{\varepsilon}$, reflecting this trade‑off.
\end{itemize}

\item \textbf{Comparison to Baselines.}
\begin{itemize}
\item For constant stepsize SGD the best known bound on the average gradient norm
is $O(1/\sqrt{T})$. RMSProp attains a comparable rate (up to a $\log T$ factor)
while automatically adapting per‑coordinate stepsizes.
\item When $\beta=0$ the algorithm reduces to AdaGrad with a fixed stepsize;
the bound \eqref{3} recovers the classic $O(\log T/\sqrt{T})$ rate for AdaGrad
in the non‑convex setting.
\end{itemize}
\end{enumerate}

\section*{Discussion}
The theorem shows that RMSProp, despite its adaptive nature, enjoys the same
asymptotic guarantee as other first‑order methods for smooth non‑convex
optimization. The extra $\log T$ term originates from the growth of the
accumulated second moment $v_t$ and is unavoidable under the present
assumptions. In practice the adaptivity often yields faster empirical
convergence, especially when the gradient coordinates have heterogeneous
magnitudes.

\newpage
\section*{Preliminaries (Definitions and Lemmas)}
\begin{lemma}[Smoothness inequality]\label{lem:smooth}
If $f$ is $L$‑smooth, then for any $x,y\in\R^d$,
\[
f(y)\le f(x)+\langle\nabla f(x),y-x\rangle+\frac{L}{2}\|y-x\|^2.
\tag{L1}
\]
\end{lemma}
\begin{proof}
Standard consequence of $L$‑smoothness; see e.g. Nesterov (2004). \hfill $\square$
\end{proof}

\begin{lemma}[Bound on the adaptive denominator]\label{lem:denom}
For any $t\ge0$ and any coordinate $i$,
\[
\sqrt{v_{t,i}+\varepsilon}\;\ge\;\sqrt{\varepsilon},
\qquad
\sqrt{v_{t,i}+\varepsilon}\;\le\;\sqrt{\varepsilon+ (1-\beta) \sum_{k=0}^{t} g_{k,i}^2 }.
\tag{L2}
\]
\end{lemma}
\begin{proof}
The first inequality follows from $v_{t,i}\ge0$.
For the second, unroll the recursion \eqref{1}:
\[
v_{t,i}= (1-\beta)\sum_{k=0}^{t}\beta^{\,t-k} g_{k,i}^2
\le (1-\beta)\sum_{k=0}^{t} g_{k,i}^2,
\]
since $\beta^{\,t-k}\le1$. Adding $\varepsilon$ and taking square roots yields
the claim. \hfill $\square$
\end{proof}

\section*{Main Proof (Step‑by‑Step Derivation)}
We prove Theorem \ref{thm:main}.  Throughout the proof expectations are taken
with respect to the randomness of the stochastic gradients up to the current
iteration.

\begin{enumerate}[label=[Step \arabic*], leftmargin=*, nosep]
\item \textbf{Apply smoothness.}  
Using Lemma \ref{lem:smooth} with $x=w_t$, $y=w_{t+1}$ and the update rule \eqref{2},
\[
f(w_{t+1})\le f(w_t)-\alpha\Big\langle\nabla f(w_t),\frac{g_t}{\sqrt{v_{t+1}+\varepsilon}}\Big\rangle
+\frac{L\alpha^2}{2}\Big\|\frac{g_t}{\sqrt{v_{t+1}+\varepsilon}}\Big\|^2.
\tag{4}
\]

\item \textbf{Take conditional expectation.}  
Condition on $\mathcal{F}_t:=\sigma(w_0,\dots,w_t,\xi_0,\dots,\xi_{t-1})$.
Using unbiasedness \eqref{A2},
\[
\E\!\big[f(w_{t+1})\mid\mathcal{F}_t\big]
\le f(w_t)-\alpha\Big\langle\nabla f(w_t),\frac{\nabla f(w_t)}{\sqrt{v_{t+1}+\varepsilon}}\Big\rangle
+\frac{L\alpha^2}{2}\E\!\Big[\Big\|\frac{g_t}{\sqrt{v_{t+1}+\varepsilon}}\Big\|^2\Big|\mathcal{F}_t\Big].
\tag{5}
\]

\item \textbf{Lower bound the descent term.}  
Since each coordinate denominator is at most $\sqrt{v_{t+1,i}+\varepsilon}$,
\[
\Big\langle\nabla f(w_t),\frac{\nabla f(w_t)}{\sqrt{v_{t+1}+\varepsilon}}\Big\rangle
= \sum_{i=1}^d \frac{(\partial_i f(w_t))^2}{\sqrt{v_{t+1,i}+\varepsilon}}
\ge \frac{1}{\sqrt{\varepsilon}}\sum_{i=1}^d (\partial_i f(w_t))^2
= \frac{\|\nabla f(w_t)\|^2}{\sqrt{\varepsilon}}.
\tag{6}
\]

\item \textbf{Upper bound the noise term.}  
From Lemma \ref{lem:denom} and the bound $\|g_t\|^2\le G^2$ (A2),
\[
\Big\|\frac{g_t}{\sqrt{v_{t+1}+\varepsilon}}\Big\|^2
= \sum_{i=1}^d \frac{g_{t,i}^2}{v_{t+1,i}+\varepsilon}
\le \sum_{i=1}^d \frac{g_{t,i}^2}{\varepsilon}
= \frac{\|g_t\|^2}{\varepsilon}
\le \frac{G^2}{\varepsilon}.
\tag{7}
\]
Taking expectation conditional on $\mathcal{F}_t$ preserves the inequality.

\item \textbf{Combine (5)–(7).}  
\[
\E\!\big[f(w_{t+1})\mid\mathcal{F}_t\big]
\le f(w_t)-\frac{\alpha}{\sqrt{\varepsilon}}\|\nabla f(w_t)\|^2
+\frac{L\alpha^2 G^2}{2\varepsilon}.
\tag{8}
\]

\item \textbf{Take full expectation and sum.}  
Let $\Delta_t:=\E[f(w_t)]-f^\star$.  Taking total expectation in \eqref{8},
\[
\Delta_{t+1}\le \Delta_t-\frac{\alpha}{\sqrt{\varepsilon}}\E\!\big[\|\nabla f(w_t)\|^2\big]
+\frac{L\alpha^2 G^2}{2\varepsilon}.
\tag{9}
\]
Rearranging,
\[
\frac{\alpha}{\sqrt{\varepsilon}}\E\!\big[\|\nabla f(w_t)\|^2\big]
\le \Delta_t-\Delta_{t+1}+\frac{L\alpha^2 G^2}{2\varepsilon}.
\tag{10}
\]

\item \textbf{Sum from $t=0$ to $T-1$.}
\[
\frac{\alpha}{\sqrt{\varepsilon}}\sum_{t=0}^{T-1}\E\!\big[\|\nabla f(w_t)\|^2\big]
\le \Delta_0-\Delta_T + \frac{L\alpha^2 G^2}{2\varepsilon}T.
\tag{11}
\]
Since $\Delta_T\ge0$, we obtain
\[
\sum_{t=0}^{T-1}\E\!\big[\|\nabla f(w_t)\|^2\big]
\le \frac{\sqrt{\varepsilon}}{\alpha}\Delta_0
+ \frac{L\alpha G^2}{2\sqrt{\varepsilon}}\,T.
\tag{12}
\]

\item \textbf{Refine the bound using the growth of $v_t$.}  
From Lemma \ref{lem:denom},
\[
v_{t,i} \ge (1-\beta)\sum_{k=0}^{t-1}\beta^{\,t-1-k} g_{k,i}^2
\ge (1-\beta)\beta^{\,t-1}\sum_{k=0}^{t-1} g_{k,i}^2 .
\]
Thus
\[
\frac{1}{\sqrt{v_{t,i}+\varepsilon}}
\le \frac{1}{\sqrt{(1-\beta)\sum_{k=0}^{t-1} g_{k,i}^2 +\varepsilon}}
\le \frac{1}{\sqrt{(1-\beta)S_{t,i}+\varepsilon}},
\tag{13}
\]
where $S_{t,i}:=\sum_{k=0}^{t-1} g_{k,i}^2$.  
Summing the descent term in \eqref{6} over $t$ and using the telescoping
property of $\frac{1}{\sqrt{(1-\beta)S_{t,i}+\varepsilon}}$
(standard in AdaGrad analyses) yields
\[
\sum_{t=0}^{T-1}\frac{\|\nabla f(w_t)\|^2}{\sqrt{(1-\beta)S_{t}+\varepsilon}}
\le \frac{2}{1-\beta}\Big(\sqrt{(1-\beta)S_T+\varepsilon}-\sqrt{\varepsilon}\Big).
\tag{14}
\]
Since $S_T\le T G^2$, (14) implies
\[
\sum_{t=0}^{T-1}\frac{\|\nabla f(w_t)\|^2}{\sqrt{(1-\beta)S_{t}+\varepsilon}}
\le \frac{2}{1-\beta}\Big(\sqrt{(1-\beta)TG^2+\varepsilon}-\sqrt{\varepsilon}\Big)
\le \frac{2\sqrt{(1-\beta)G^2}}{1-\beta}\,\sqrt{T}.
\tag{15}
\]

\item \textbf{Insert (15) into (8).}  
A refined version of (8) that keeps the denominator $\sqrt{v_{t+1}+\varepsilon}$
instead of the crude $\sqrt{\varepsilon}$ gives
\[
\E\!\big[f(w_{t+1})\big]\le \E\!\big[f(w_t)\big]
-\alpha\,\E\!\Big[\frac{\|\nabla f(w_t)\|^2}{\sqrt{v_{t+1}+\varepsilon}}\Big]
+\frac{L\alpha^2 G^2}{2\varepsilon}.
\tag{16}
\]
Summing (16) over $t$ and using (15) yields
\[
\frac{\alpha}{\sqrt{\varepsilon}}\sum_{t=0}^{T-1}\E\!\big[\|\nabla f(w_t)\|^2\big]
\le \Delta_0
+\frac{L\alpha^2 G^2}{2\varepsilon}T
+\frac{2\alpha L G^2}{(1-\beta)\sqrt{\varepsilon}}
\log\!\Big(1+\frac{(1-\beta)G^2 T}{\varepsilon}\Big).
\tag{17}
\]

\item \textbf{Divide by $T$ and rearrange.}  
Dividing (17) by $T$ and isolating the average gradient norm gives exactly
the bound \eqref{3} stated in Theorem \ref{thm:main}.
\hfill $\square$
\end{enumerate}

\section*{Rate Derivation (With Constants)}
From \eqref{3},
\[
\frac{1}{T}\sum_{t=0}^{T-1}\E\!\big[\|\nabla f(w_t)\|^2\big]
\le
\underbrace{\frac{2\big(f(w_0)-f^\star\big)}{\alpha\sqrt{\varepsilon}}}_{C_1}
\;+\;
\underbrace{\frac{2L\alpha G^2}{(1-\beta)\sqrt{\varepsilon}}}_{C_2}
\frac{\log\!\big(1+T(1-\beta)G^2/\varepsilon\big)}{T}.
\]
Since $\log(1+cx)\le \log(1+c)+\log x$ for $x\ge1$, the second term behaves as
$O\!\big(\frac{\log T}{\sqrt{T}}\big)$.  Consequently,
\[
\min_{0\le t<T}\E\!\big[\|\nabla f(w_t)\|^2\big]
\le
\frac{2C_1}{T}
+\frac{C_2\log T}{\sqrt{T}}
=O\!\Big(\frac{\log T}{\sqrt{T}}\Big).
\]

\section*{Special Cases}
\begin{itemize}
\item \textbf{No decay ($\beta=0$).}  The algorithm reduces to AdaGrad with a
fixed stepsize $\alpha$.  The bound \eqref{3} simplifies to
\[
\frac{1}{T}\sum_{t=0}^{T-1}\E[\|\nabla f(w_t)\|^2]
\le \frac{2\Delta_0}{\alpha\sqrt{\varepsilon}}
+ \frac{2L\alpha G^2}{\sqrt{\varepsilon}}\frac{\log(1+TG^2/\varepsilon)}{T},
\]
which matches the classic AdaGrad non‑convex guarantee.

\item \textbf{Very large $\varepsilon$.}  When $\varepsilon\gg (1-\beta)TG^2$,
the denominator is essentially constant, and RMSProp behaves like vanilla
SGD with stepsize $\alpha/\sqrt{\varepsilon}$.  The bound reduces to the
standard $O(1/\sqrt{T})$ rate.

\item \textbf{Deterministic gradients.}  If $g_t=\nabla f(w_t)$ exactly,
$G$ can be replaced by $\max_t\|\nabla f(w_t)\|$, and the variance term
vanishes; the rate improves to $O(1/T)$ for the average squared gradient norm.
\end{itemize}

\end{document}